% =========================================================
% Proof Stub: Cauchy Criterion (Cauchy iff Convergent in R)
% Source: volume-iii/analysis/sequences/notes/sequences/notes-cauchy.tex
% Mode: standard
% =========================================================

\clearpage
\phantomsection
\hypertarget{proof-RA-ABB-P-cauchy-criterion}{}

\subsubsection[Cauchy Criterion in $\mathbb{R}$]{Proof --- Cauchy Criterion in $\mathbb{R}$}

\bigskip

\noindent
\textbf{Source.}
Abbott, \emph{Understanding Analysis}, Theorem~2.6.4; Tao, \emph{Analysis I}.
See notes \texttt{notes-cauchy.tex}.

\vspace{0.75em}

\noindent
\textbf{Goal.}
A real sequence converges if and only if it is Cauchy.

\vspace{0.75em}

\noindent
\textbf{Logical form.}
\[
  (x_n) \text{ converges in } \mathbb{R}
  \;\iff\;
  (x_n) \text{ is Cauchy.}
\]

\vspace{0.75em}

\noindent
\textbf{Proof strategy.}
$(\Rightarrow)$: Already proved (Convergent implies Cauchy).
$(\Leftarrow)$: Cauchy implies bounded (proved), so by Bolzano-Weierstrass
$(x_n)$ has a convergent subsequence $(x_{n_k}) \to L$.
Show $(x_n) \to L$: given $\varepsilon > 0$, use the Cauchy condition to get
$N$, and the subsequence convergence to get $k_0$ with $n_{k_0} \ge N$ and
$|x_{n_{k_0}} - L| < \varepsilon/2$.
Then for $n \ge N$: $|x_n - L| \le |x_n - x_{n_{k_0}}| + |x_{n_{k_0}} - L| < \varepsilon$.

\vspace{0.75em}

\noindent
\textbf{Note.}
The $(\Leftarrow)$ direction is one of the main payoffs of this chapter.
It weaves together Cauchy boundedness, Bolzano-Weierstrass, and subsequence
inheritance. Take time to map the proof dependencies carefully.

\vspace{0.75em}

\noindent
\textbf{Proof.}
\begin{proof}
\Claim $(x_n)$ converges in $\mathbb{R}$ iff $(x_n)$ is Cauchy.

\Given $(x_n)$ a real sequence.

\Goal Prove both directions.

\Strategy $(\Rightarrow)$: refer to the previously proved direction.
$(\Leftarrow)$: chain Cauchy-bounded, BW, subsequence-inherits-limits, Cauchy.

\medskip

\textbf{($\Rightarrow$) Convergent $\Rightarrow$ Cauchy:}
% --- previously proved ---

\medskip

\textbf{($\Leftarrow$) Cauchy $\Rightarrow$ Convergent:}
% --- (x_n) Cauchy ⇒ (x_n) bounded (Cauchy sequences are bounded) ---

% --- by BW: extract convergent subsequence (x_{n_k}) → L ---

% --- given ε > 0: ---
%   use Cauchy: get N with |x_m - x_n| < ε/2 for m,n ≥ N ---
%   find k₀ with n_{k₀} ≥ N and |x_{n_{k₀}} - L| < ε/2 ---

% --- for n ≥ N: ---
%   |x_n - L| ≤ |x_n - x_{n_{k₀}}| + |x_{n_{k₀}} - L| < ε/2 + ε/2 = ε ---

\end{proof}

\begin{remark*}[Proof shape]
The hard direction $(\Leftarrow)$ is a three-ingredient argument:
Cauchy$\Rightarrow$bounded $\Rightarrow$ (via BW) convergent subsequence
$\Rightarrow$ (via Cauchy condition) full sequence converges.
\end{remark*}

\begin{remark*}[Dependencies]
\begin{itemize}
  \item Convergent implies Cauchy.
  \item Cauchy sequences are bounded.
  \item Bolzano-Weierstrass Theorem.
  \item Subsequences inherit limits.
\end{itemize}
\end{remark*}
